\originalchapter{活动星系核的尺度、能源与吸积}

本章对应 8.902 第9讲的 AGN 部分. 活动星系核是星系中心显著的非恒星辐射源, 并非要求它在每个波段都压倒整个宿主. 要判断中央能源, 可以同时利用亮度、变化时间、谱线速度和辐射机制. 这条观测证据链比单凭一张有喷流的图片更具体; 本卷不复制原讲义中带许可例外的照片.

\originalsection{分类反映辐射与观测方向}
射电星系的喷流可把能量输送到远处射电瓣, 相对论电子的同步辐射形成明显非热谱. 类星体在光学呈明亮且紧致的核; 并非所有光学类星体都有很强射电喷流. 历史术语 quasar 与 QSO 的射电选择背景不同, 当代广义分类不能只按名称断言有无喷流. Seyfert 核通常在可见宿主中较易辨认, 核光度比最亮类星体低. BL Lac 等 blazar 的喷流接近视线, 相对论束射使连续谱增强、变化快且可有较高偏振.

课程所列类星体总光度约 $10^{45}$--$10^{48}\,\mathrm{erg\,s^{-1}}$, Seyfert 核约 $10^{42}$--$10^{45}\,\mathrm{erg\,s^{-1}}$, 是样本数量级而非严格分类边界. 原课射电项将 $10^{18}$ 的阈值标作 $\Ls$, 与随后的总能源尺度难以自洽; 该单位未能核定, 本卷不以它定义射电强弱或计算黑洞质量. 观测射电强度需说明频率、光谱光度密度和选择阈值, 不能擅自把原数改配一个猜测单位.

\originalsection{变化时间限制发光区大小}
若区域的相干变化受内部信息传播控制且无强束射, 本征时间窗 $\Delta t$ 给发光区尺度 $R\lesssim c\Delta t$. 对红移源, $\Delta t_{\rm obs}=(1+z)\Delta t$. 若喷流有 Doppler 因子 $\mathcal D$, 则本征区尺度上界近似
\[
R'\lesssim\frac{c\Delta t_{\rm obs}\mathcal D}{1+z},\qquad
\mathcal D=[\Gamma(1-\beta\cos i)]^{-1}.
\]
因而不能把所有快变对象都按静止源的 $c\Delta t_{\rm obs}$ 估大小. 无红移与束射时, 一天光程约 $8.4\times10^{-4}\,\mathrm{pc}$, 与课程给出的日尺度核尺寸相符.

黑洞 Schwarzschild 半径为 $r_s=2GM_\bullet/c^2\simeq9.57\times10^{-14}(M_\bullet/\Ms)\,\mathrm{pc}$. 因此百万太阳质量给 $r_s\sim10^{-7}\,\mathrm{pc}$, 十亿太阳质量给 $10^{-4}\,\mathrm{pc}$. 快变区可以在黑洞附近, 但大小上界本身既不等于视界半径, 也未唯一决定质量. 宽线气体的速度和响应时间可再提供质量尺度 $M_\bullet\sim f_{\rm geom}R_{\rm BLR}\Delta v^2/G$, 其中 $R_{\rm BLR}\sim c\tau_{\rm rest}$, 几何和投影因子需另标定.

高亮度集中在小区域, 使单靠普通恒星密集堆积或频繁超新星供能的模型面临质量、寿命和密度约束. 例如 $L=10^{46}\,\mathrm{erg\,s^{-1}}$ 一年释放 $3.16\times10^{53}\,\mathrm{erg}$, 若每次事件可辐射 $10^{51}\,\mathrm{erg}$, 要数百次每年且集中在极小核区. 这是数量级排查, 未把所有瞬变核都当作被同一单数值严格排除的模型. 对持续高功率活动核, 引力吸积提供更有效的能源通道.

\originalsection{吸积效率与薄盘加热}
气体落入深引力势可把部分静质量能转成辐射, 写 $L=\eta\dot M_{
\rm in}c^2$. $\eta$ 是辐射效率, 随吸积模式、内边界、自旋和能量损失改变, 不是所有黑洞同一常数. 课程取 $\eta\simeq1/16$ 示范, 与氢燃烧的约0.007静质量转化率比较; 原始核燃料中实际可燃烧比例又是不同因素.

气体有角动量时不会直接径向落下. 在薄、近 Kepler、稳态且内缘零力矩模型中, 角动量向外输运使气体向内流. 以运动黏性系数 $\nu_{\rm visc}$ 表示应力, 角动量守恒给
\[
\nu_{\rm visc}\Sigma=\frac{\dot M_{
\rm in}}{3\pi}\left[1-\sqrt{\frac{r_{
\rm in}}r}\right].
\]
每面局部耗散通量为 $F=\tfrac12\nu_{\rm visc}\Sigma(r\Omega')^2$. 用 $\Omega^2=GM_\bullet/r^3$, 得
\[
F(r)=\frac{3GM_\bullet\dot M_{
\rm in}}{8\pi r^3}
\left[1-\sqrt{\frac{r_{
\rm in}}r}\right].
\]
这说明热量来自剪切应力与引力能源的共同预算. 若光厚且局部辐射有效, $\sigma_{\rm SB}T_{
\rm eff}^4=F$, 远离内缘 $T_{
\rm eff}\propto r^{-3/4}$. 不同半径的热谱叠加形成多温度盘; X 射线可来自更热的日冕等额外成分, 不应把课程给出的 $10^8$ K 无条件当作所有超大质量黑洞薄盘的有效温度.

双面面积积分给牛顿薄盘 $L=\int_{r_{
\rm in}}^\infty4\pi rF(r)\dd r=GM_\bullet\dot M_{
\rm in}/(2r_{
\rm in})$. 该近似靠近视界时失效; Schwarzschild 的相对论薄盘效率约0.057, 来自 ISCO 结合能, 见已有 GR 卷前置结果. 本章不将 Newton 近似在 $r_{
\rm in}=6GM/c^2$ 的代入值冒充完整 GR 计算, 也未从一个有效黏性系数求出真实磁湍流机制.

\originalsection{Eddington 光度的假设与增长时间}
对近球对称辐射通量 $L/(4\pi r^2)$, 单位质量辐射力 $a_{
\rm rad}=\kappa_F L/(4\pi r^2c)$, 其中 $\kappa_F$ 是按实际光谱加权的不透明度. 与引力 $GM_\bullet/r^2$ 比较, 得
\[
L_{
\rm Edd}=\frac{4\pi GM_\bullet c}{\kappa_F}.
\]
完全电离氢且以 Thomson 散射为主时 $\kappa_F=\sigma_T/m_p$, 给 $L_{
\rm Edd}\simeq1.26\times10^{38}(M_\bullet/\Ms)\,\mathrm{erg\,s^{-1}}$. 电子受辐射力而通过等离子体耦合拖动质子, 所以配的是每质子质量, 不是直接把 $m_e$ 作为被引力束缚的全部质量.

称为“上限”时需保留几何、稳态、光子俘获和不透明度条件. 非球对称盘、束射或超临界流可使表观光度或质量流量与这份简单平衡不同, 黑洞也不会因辐射压力就自身被吹散. 定义给定效率的参考流入率 $\dot M_{
\rm Edd}=L_{
\rm Edd}/(\eta c^2)$,
\[
\dot M_{
\rm Edd}\simeq
2.22\times10^{-8}\left(\frac{0.1}{\eta}\right)
\frac{M_\bullet}{\Ms}\,\Ms\,\mathrm{yr^{-1}}.
\]
若 $\eta=1/16$, 系数约 $3.55\times10^{-8}$, 与原课的 $5\times10^{-10}$ 不相符; 此处从光度、$c^2$ 与秒到年的单位换算重新确定, 不保留错误系数.

若辐射拿走质量能的比例 $\eta$, 黑洞增长率为 $\dot M_\bullet=(1-\eta)\dot M_{
\rm in}$. 固定 Eddington 比 $\lambda=L/L_{
\rm Edd}$ 下, 质量指数增长, e 折叠时间
\[
t_{
\rm grow}=\frac{\eta}{(1-\eta)\lambda}
\frac{\kappa_Fc}{4\pi G}.
\]
对上述 Thomson 氢标定, 后因子约 $4.5\times10^8$ 年. $\eta=0.1$, $\lambda=1$ 给约五千万年. 这个模型比较高红移大黑洞形成所需时间时, 还应说明种子质量、间歇活动和效率史, 不将一次指数积分当作已唯一解释所有早期黑洞.

\originalsection{同步辐射与统一结构的用途}
Gaussian 单位下, 在磁场 $B$ 中运动的相对论电子, 典型同步频率为
\[
\nu_{
\rm syn}\sim\frac{3eB\gamma^2\sin\alpha}{4\pi m_ec},
\]
其中 $\alpha$ 为俯仰角, $\gamma$ 为电子 Lorentz 因子. 能量越高, 辐射频率越高, 冷却也更快. 对等向俯仰角平均, 同步损失功率 $P\simeq(4/3)\sigma_TcU_B\gamma^2\beta^2$, $U_B=B^2/(8\pi)$, 给 $t_{
\rm cool}\sim\gamma m_ec^2/P$. 非热谱、偏振和喷流方向可以约束这些量, 但仅由某个频率不能单独求出 $B$ 与 $\gamma$, 因为它们以乘积进入.

AGN 的简化统一图包含中央黑洞与盘、近核宽线区、较远窄线区、尘埃遮挡结构及可能的喷流. 宽线区轨道速度较高, 视线速度分布使谱线更宽; 窄线区较慢, 名称并非单纯由气体“厚”或“薄”决定. 尘埃吸收光并辐射红外, 遮挡方向使某些核的宽线难见. 喷流朝视线时束射增强, 对应 blazar 的一类表现.

\begin{center}
\begin{tikzpicture}[>=Latex]
\fill(0,0)circle(2pt)node[below]{黑洞};
\draw[structurecolor,thick](-0.6,0.08)--(0.6,0.08)node[right]{盘};
\draw[dashed](-1,-0.4)ellipse[x radius=.3,y radius=.5];
\draw[dashed](1,-0.4)ellipse[x radius=.3,y radius=.5];
\node at(1.75,-0.65){尘埃};
\draw[->,thick](0,.2)--(0,2)node[above]{喷流};
\draw[->](2.6,1.6)--(.25,.35)node[midway,above]{观测方向};
\node at(-1.4,1.1){宽线区};\node at(2.4,.65){窄线区};
\end{tikzpicture}
\end{center}
图只表示组成与方向, 不按长度比例绘制, 不替代某个真实活动核的几何测量. 统一模型解释许多方向相关的分类, 仍须允许不同吸积率、自旋和环境. AGN 的风与喷流还能向宿主输入能量, 改变前章气体流入流出与恒星形成; 反馈效能需独立能量预算, 不是仅看到黑洞就自动确定每个星系的化学史.
